\cotefile{63}
\scanpage{21}\mspage{6}
\[
\mathrm{u}^{4}\,\mathrm{c}'_{4}=\mathrm{c}_{4}\qquad
\mathrm{u}^{6}\,\mathrm{c}'_{6}=\mathrm{c}_{6}\qquad
\mathrm{u}^{12}\,\Delta'=\Delta\qquad.
\tag*{(1.8)}
\]

Les formules (1.8) expriment que la section $\mathrm{c}_{4}\,\pi^{\handwritten{\otimes}4}$ (resp. $\mathrm{c}_{6}\,\pi^{\handwritten{\otimes}6}$, resp. $\Delta\,\pi^{\handwritten{\otimes}12}$) du fibré $\omega^{\handwritten{\otimes}4}$ (resp. $\omega^{\handwritten{\otimes}6}$, $\omega^{\handwritten{\otimes}12}$)\raisebox{.35ex}{\tiny.} ne dépend pas du choix arbitraire de $\pi$, $\mathrm{x}$ et $\mathrm{y}$.
\par\noindent\hfill
\margin{\hspace{1.8cm}\makebox[0pt][r]{\begin{minipage}[t]{10.1cm}
\diagram{\handwritten{\raisebox{2.65\baselineskip}{\uline{Sur un corps}}\ \(\left\{\begin{array}{@{}l@{\quad}l@{}}
\text{Cas elliptique :}&\Delta\neq0\\
\text{cas multiplicatif :}&\Delta=0,\ c_{4}\neq0\ \text{(équiv.}\\[-.2ex]
&\hfill{\scriptstyle\unsure{c|e}}\\[-.3ex]
&\hspace{4em}\text{à }c_{6}\neq0)\\
\text{cas additif :}&\Delta=0,\ c_{4}=0\ \text{i.e.}\\
&\Delta=0,\ c_{6}=0
\end{array}\right.\)}}
\end{minipage}}}
\par\vspace{7.7\baselineskip}

2.\qquad \uline{Cas où 2 et 3 sont inversibles}.\margin{\handwritten{$\backslash$}}

Lorsque $2$ et $3$ sont inversibles, on peut d'une et d'une seule façon choisir $\mathrm{x}$ et $\mathrm{y}$ ($\pi$ étant donné) de sorte que dans l'équation (1.3.) on ait $\mathrm{a}_{1}=\mathrm{a}_{2}=\mathrm{a}_{3}=0$. On a alors
\[
\mathrm{y}^{2}=\mathrm{x}^{3}+\mathrm{a}_{4}\,\mathrm{x}+\mathrm{a}_{6}\qquad,\ \text{avec}\quad\pi=\mathrm{dx}/2\mathrm{y}\qquad,
\tag*{(2.1)}
\]
\[
\mathrm{c}_{4}=-2^{4}.3\mathrm{a}_{4}\quad,\quad
\mathrm{c}_{6}=-2^{5}.3^{3}\,\mathrm{a}_{6}\quad,\quad
\Delta=2^{4}(-4\mathrm{a}_{4}^{3}-27\,\mathrm{a}_{6}^{2})\quad.
\tag*{(2.2)}
\]

Posant $\mathrm{Y}=2\mathrm{y}$, $\mathrm{X}=\mathrm{x}$, on retombe sur la forme de Weierstrass
\[
\mathrm{Y}^{2}=4\,\mathrm{X}^{3}-\mathrm{g}_{2}\,\mathrm{X}-\mathrm{g}_{3}\quad,\ \text{avec}\quad\pi=\mathrm{d}\,\mathrm{X}/\mathrm{Y}\quad,\quad\text{et}
\tag*{(2.3)}
\]
\[
\begin{aligned}[t]
&\mathrm{g}_{2}=-4\,\mathrm{a}_{4}=\handwritten{\frac}{1}{12}\,\mathrm{c}_{4}\qquad
\mathrm{g}_{3}=-4\,\mathrm{a}_{6}=\handwritten{\frac}{1}{216}\,\mathrm{c}_{6}\\
&\Delta=\mathrm{g}_{2}^{3}-27\,\mathrm{g}_{3}^{2}\qquad.
\end{aligned}
\tag*{(2.4)}
\]

\uline{Proposition} 2.5.\quad \uline{Au dessus de} $\operatorname{Spec}(\handwritten{\mathbb{Z}}[2^{-1},\,3^{-1}]$, \uline{il existe un schéma de modules pour les courbes de genre un munies d'une forme différentielle invariante inversible}. \uline{C; schéma est}
\[
\operatorname{Spec}(\handwritten{\mathbb{Z}}[2^{-1},\,3^{-1}][\mathrm{g}_{2},\,\mathrm{g}_{3}]\qquad,
\]
\clearpage
% archive sheet number: 20
